% Generated draft: requires mathematical and source review.
\documentclass[11pt]{article}
\usepackage[margin=25mm]{geometry}
\usepackage{fontspec}
\setmainfont{Latin Modern Roman}
\usepackage{amsmath,amssymb,mathtools,unicode-math}
\setmathfont{Latin Modern Math}
\usepackage{enumitem,needspace,xurl,hyperref}
\setlength{\parindent}{0pt}
\setlength{\parskip}{5pt}
\setlength{\emergencystretch}{3em}
\newcommand{\N}{\mathbb{N}}
\newcommand{\Z}{\mathbb{Z}}
\newcommand{\Q}{\mathbb{Q}}
\newcommand{\R}{\mathbb{R}}
\newcommand{\C}{\mathbb{C}}
\newcommand{\F}{\mathbb{F}}
\newcommand{\A}{\mathbb{A}}
\newcommand{\PP}{\mathbb{P}}
\newcommand{\E}{\mathbb{E}}
\newcommand{\eps}{\varepsilon}
\newcommand{\dd}{\,\mathrm{d}}
\begin{document}
\section*{UnsolvedMath ID: 30003484}
\textbf{Source problem code:} catalogue entry
\subsection{Definitions and mathematical statement}
For a fixed leaf number \(n\), consider two distinct \(n\)-leaf phylogenetic networks \(N_1\) and \(N_2\). The source does not specify the admissible network class, the equivalence relation determining distinctness, or whether leaf labels must agree. These choices must be fixed before the problem has a fully specified domain.

For each network \(N\), write \(M_N\) for the associated model set, called the model variety in the supplied record, and put
\[
V_N=\overline{M_N}^{\mathrm{Zar}}.
\]
This notation requires a choice of model, ground field, and ambient algebraic space. For the following intersection to be meaningful, the two model sets must be interpreted in a common ambient space \(X\), with compatible coordinates. These are prerequisites for the formulation, not specifications supplied by the source.

\textbf{Problem.} With those choices fixed, characterize the pairs of distinct \(n\)-leaf phylogenetic networks \((N_1,N_2)\) for which
\[
V_{N_1}\cap V_{N_2}\subsetneq V_{N_1}
\qquad\text{and}\qquad
V_{N_1}\cap V_{N_2}\subsetneq V_{N_2}.
\]
Here \(\subsetneq\) means strict set inclusion; the intersection has the induced Zariski-closed structure. The requested conditions concern properness in both varieties. The source does not specify whether an empty intersection is permitted as a proper subvariety, whether irreducibility is assumed, or whether a stronger dimension condition is intended. No expected-dimension or transversality condition is imposed in this draft.

The supplied record labels the problem partially solved and reports generic distinguishability for important large-cycle and level-one network classes, while stating that proper intersections for arbitrary network varieties remain incompletely characterized. This is supplied dated curation, not an independently verified mathematical or literature assessment.
\subsection{Short English statement}
Characterize when the model varieties of two distinct phylogenetic networks with the same number of leaves intersect in a strictly smaller closed subset of each. A precise instance requires specifying the network class, statistical model, ground field, and common ambient coordinates; the supplied record leaves these choices unresolved.
\subsection{Sources}
\begin{itemize}
\item Supplied local UnsolvedMath registry record, permanent ID \(30003484\), title \emph{Proper Intersections of Phylogenetic Network Varieties}, source problem code \texttt{catalogue entry}, in the Oberwolfach Reports Open Problems collection. The record gives proposed year \(2017\) and names no proposer.
\item Supplied report citation: \emph{Algebraic Statistics (2018)}, Oberwolfach Reports, pp.\ 1207--1279, DOI \texttt{10.4171/owr/2017/20}. \url{https://doi.org/10.4171/owr/2017/20}. The background describes Oberwolfach Reports as documenting mathematical workshops and their open-problem discussions. The supplied year \(2018\) and the \(2017\) appearing in the DOI are retained without reconciliation.
\item Supplied literature source: ``Gross and Long on distinguishing networks.'' \url{https://doi.org/10.1137/17M1134238}.
\item Supplied literature source: ``Brits and collaborators on full level-one identifiability.'' \url{https://arxiv.org/abs/2607.12919}.
\item Supplied literature source: ``algebraic matroid identifiability.'' \url{https://doi.org/10.1016/j.jsc.2020.04.012}.
\item Literature assessment date: \texttt{2026-08-21}. The supplied dated curation labels the problem partially solved and states that generic distinguishability is proved for important large-cycle and level-one network classes, but conditions for proper intersections of arbitrary network varieties are not fully characterized. The background describes this as dated literature triage derived from the supplied OWR curation and requests independent verification of time-sensitive status and cited claims. This assessment and the references have not been independently verified here.
\item The registry assigns the statement status \texttt{corrected\_verified} and says that mathematical content was preserved while notation, typography, and sentence structure were repaired for publication. This is a supplied editorial qualification, not independent mathematical verification of this draft.
\end{itemize}
\end{document}
